% Resumo
\begin{resumo}
\vspace{-1cm}
\onehalfspacing

\noindent Este trabalho apresenta um estudo comparativo experimental de três
implementações do algoritmo Quicksort: a versão recursiva clássica, uma versão
híbrida em que a partição é interrompida para subvetores com menos de $M$ elementos,
os quais são ordenados por Insertion Sort, e uma versão híbrida com escolha de pivô
por mediana-de-três. As implementações foram desenvolvidas em C++17 sobre um motor
único parametrizado, de modo a garantir que as três variantes contabilizem
comparações e trocas de forma idêntica. Os algoritmos foram avaliados sobre cinco
massas de teste --- aleatória, ordenada, inversamente ordenada, com muitos elementos
repetidos e quase ordenada --- em quatro tamanhos de entrada, entre mil e quinhentos
mil elementos, medindo-se tempo de execução, número de comparações e número de
trocas. A calibração empírica mostrou que o mínimo do número de comparações é estável
em $M \approx 20$--$25$, enquanto o mínimo de tempo é ruidoso; adotou-se $M = 40$ pelo
critério do menor valor dentro de 2\% do melhor tempo. Esse valor é superior à
faixa de 5 a 25 recomendada pela literatura clássica, diferença atribuída às
características do hardware atual. As versões híbridas reduziram o tempo de execução
entre 16\% e 54\% em relação à versão recursiva pura. Um experimento específico
forçou o pior caso do algoritmo, combinando vetor ordenado com escolha inadequada de
pivô, e confirmou numericamente a previsão teórica: a razão entre o número de
comparações e $n^2$ convergiu para exatamente $0{,}5000$, com diferença de
aproximadamente duas mil vezes em relação à versão com mediana-de-três. A análise
crítica revelou que, em dados aleatórios, a versão híbrida executa mais comparações
que a versão recursiva e ainda assim é mais rápida, evidenciando que a contagem de
operações elementares deixou de ser preditor confiável do tempo de execução em
hardware moderno.

\vspace*{.75cm}

\noindent \textbf{Palavras-chave:} Quicksort, Insertion Sort, Algoritmos de
Ordenação, Análise de Complexidade, Análise Experimental de Algoritmos,
Mediana-de-três.

\end{resumo}
